The main lesson is that durable unlearning requires a different evidence standard than immediate forgetting.
In the local MVP, a method that looked best on benign survival curves failed once the comparison controlled for $r@0$ and used answer-bearing stress.
This is exactly the kind of failure the harness is meant to expose.

The results also suggest that the relearning pool matters.
Benign pools test stability under ordinary updates, but they are too weak to establish durable forgetting.
Answer-bearing pools are harsher and more informative because they probe whether a suppressed fact can be reintroduced through plausible post-unlearning updates.

The Qwen experiments add a second lesson.
When Strong NPO could not match HLC's immediate forgetting, constrained gradient ascent provided a useful frontier baseline.
This baseline is not SOTA in the broader literature, but it is locally decisive: if HLC cannot beat an immediate-matched constrained GA baseline, it cannot support a strong method claim in this harness.

The strongest paper angle is therefore methodological.
The harness gives reviewers concrete evidence that the authors are not mistaking aggressive forgetting for durability.
That is valuable even if the current HLC objective remains negative.
